\documentclass[11pt,dvipdfmx]{abst5lab}
\pagestyle{empty}
\usepackage{graphicx}
\usepackage{color}
\usepackage{fancyhdr}
\usepackage{amssymb}
\usepackage{amsmath}
\usepackage{array}
\usepackage{tikz}
% \usepackage[version=3]{mhchem}
\usepackage{amsfonts}
\usepackage{bm}
\usepackage{here} % 図や表の位置を固定

\title{テスト}
\etitle{TITLE}
\author{六車光貴}
\affil{京都大学生存圏研究所 橋口研究室}


%\setcounter{page}{2000}
%\pagestyle{empty}
%\lhead{}
\chead{Headerの真ん中}
%\rhead{}
\cfoot{}


\begin{document}
\maketitle

\thispagestyle{empty}

\section{研究背景と目的}


「内容」

\section{（仮）原理}

「内容」


\section{（仮）方法}
「内容」


\section{結果}
「内容」

図\ref{im:apple}にりんごの画像を示した．
画像は.png，.jpg，.pdfに対応している．

\begin{figure}[H]
  \centering
  \includegraphics[width=5cm]{fig/apple.jpg}
  \caption{林檎の図}
  \label{im:apple}
\end{figure}

\section{まとめと今後の課題}


「内容」

\begin{thebibliography}{9}
 
  \bibitem{bibiname1}
    「著者名」, ``「タイトル」, ``「掲載書籍」, 「ページ数」, 「発行年」 
  \bibitem{bibiname2}
   「著者名」, ``「タイトル」, ``「掲載書籍」, 「ページ数」, 「発行年」 
  \bibitem{bibiname3}
   「著者名」, ``「タイトル」, ``「掲載書籍」, 「ページ数」, 「発行年」 


\end{thebibliography}


\end{document}





\begin{figure}[H]
\centering
\includegraphics[height=60mm]{接触.pdf}
  \caption{整流性接触・オーム性接触のV-I特性} 
  \label{fig:接触}
\end{figure}

\begin{figure}[H]
  \centering
  \begin{minipage}{0.4\columnwidth}
  \centering
  \includegraphics[width=\columnwidth]{実験回路1.pdf}
  \caption{測定回路1} 
  \label{fig:実験回路1}
  \end{minipage}
  \begin{minipage}{0.4\columnwidth}
  \centering
   \includegraphics[width=\columnwidth]{実験回路2.pdf}
  \caption{測定回路2} 
  \label{fig:実験回路2}
  \end{minipage}
\end{figure}%


\begin{table}[H]
\caption{室温での測定結果}
\label{室温結果}
\begin{center}
\begin{tabular}{|c|c|c|c|}\hline
 試料&$\sigma$(1/$\Omega$m)& $\mu _H$($m^2 /V\cdot s$) &n($m^{-3}$)\\ \hline
試料１ &$1.19\times 10^2$ &$5.76\times 10^{-2}$&$1.52\times 10^{22}$ \\
試料２ & $6.71\times 10^2$ &$4.53\times 10^{-2}$ &$1.09\times 10^{23}$\\ \hline
\end{tabular}
\end{center}
\end{table}

{\tt }



